\section{Security Discussion and Evidence Boundaries}

The current artifact gives three kinds of evidence: end-to-end recoverability and overhead from the 10x10 grid, small aligned same-backend comparisons, and implementation-side approximation audits from the released baseline and $h_{\min}$ pilot artifacts. These evidence types answer different questions, and it is important not to blur them together.

\paragraph{Recoverability and protocol behavior.} The fixed 10x10 grid is strongest as a systems measurement. It shows that the released implementation can encode, transmit, and decode authenticated packets under one pinned local stack without decode mismatches. The aligned pilot then adds one narrower comparison point: under the same backend and prompt regime, Ghostext's segmented packet framing is more stable than the two single-segment variants we tested. Both of these are meaningful results, but neither should be rephrased as a direct detector claim. They say that the protocol works reliably and that some simpler embeddings stall more often under the same local backend. They do not say that a real censor cannot separate Ghostext text from natural text.

\paragraph{Approximation audit versus detector evidence.} The released baseline artifacts also include an implementation-side audit of the per-step approximation terms introduced in \S5. In that audit, the released $h_{\min}$ pilot yields mean total $\widehat{\mathrm{TV}}$ values of 24.50 at $h_{\min}=0.0$, 23.13 at $h_{\min}=0.5$, and 24.96 at $h_{\min}=1.0$, with corresponding packet-prefix means of 24.09, 22.63, and 24.14. These numbers are useful because they turn hidden implementation effects into explicit quantities. But they are still internal audit aggregates, not end-to-end detector outcomes, KL estimates over full continuations, or indistinguishability guarantees. A detector may exploit cross-token dependencies, prompt leakage, length cues, or backend-specific artifacts that such stepwise accounting does not capture.

\paragraph{Why prompt knowledge still matters.} All released experiments should be read as template-known or prompt-family-known operational floors. The prompt set is public, the runtime is pinned, and the evaluation intentionally privileges reproducibility over adversarial variety. This is the right choice for an implementation paper, but it means that stronger censor models remain open. A prompt-known censor could try to compare Ghostext outputs against ordinary generations under the same prompt family. A multi-message censor could exploit repeated use of one prompt family, one backend, or one sender policy. Neither setting is addressed by the present artifact.

\paragraph{Length and traffic-side surfaces.} The current artifact also does not hide all side information around a covert exchange. Packet size influences text length, retry count influences latency, and prompt family constrains stylistic range. In the released evaluation we record these quantities on purpose because they are part of reproducibility. From a security perspective, however, they are reminders that a realistic censor might use more than token-level text statistics. Padding, cover-source management, traffic shaping, or multi-message scheduling would all belong to a broader channel-level defense story, and none of those countermeasures are implemented here.

\paragraph{Further study.} A more ambitious detector-facing study would need at least four additions. First, direct comparisons against natural-cover samples under the same prompt families and model backend. Second, stronger detector baselines and modern LLM-era steganalysis pipelines. Third, cross-hardware or cross-build replay experiments to separate algorithmic failure from backend drift. Fourth, distributional measurements that connect the protocol's stepwise audit terms to empirical text-level divergence. Until those are added, the safest reading of Ghostext is as a reproducible, fail-closed covert channel prototype with explicit systems-level evidence, not as a completed detector-resistance claim.
